\section{The constrained local Gram test}
\label{sec:local-gram}

\status{This section reconstructs the projection--derivative argument behind CL9. It uses the explicitly identified signed-path basis, integral-projector and factorial-unit lemmas of the companion IP, IW and DT notes. The construction gives the source constants $(26,2,24,60)$ and, more strongly, permits the remainder $R=0$. The expanded note proves the intermediate assertions and records their exact source dependencies.}

\subsection{The conditions and the two orders}

Work over the source DVR with fraction field $K$, prime $p\ge5$, and $v(p)=1$; arbitrary ramification is allowed. Each letter has $r$ coordinates and is a vector or covector. Write
\[
 \langle F,H\rangle=\sum_\alpha F_\alpha H_\alpha\alpha!,\quad
 b(F)=\min_{F_\alpha\ne0}\left(v(F_\alpha)+\tfrac12v_p(\alpha!)\right),\quad
 n_0=\tfrac12v_p\!\left(\prod_xd_x!\right).
\]
For a set $S$ of letters, $\operatorname{ord}_S F$ is the least exponent of $t$ after applying a generic conjugate of $\operatorname{diag}(1,\ldots,1,t)$ to those letters, with the dual action on covectors. It is distinct from the coefficient valuation $v$.

Segments $j\in\mathbb Z/J\mathbb Z$ contain signed base letters $L_j$ and vector letters $z_{j,m}$, $0\le m<r$. Put $P_{j,m}=p^{g_{j,m}}$. The enlarged polynomial ring replaces each $z$ by a label $a$, carry $u$, and tangent $\tau$, with
\[
 \phi(a_{j,m,c})=z_{j,m,c},\qquad
 \phi(u_{j,m,c})=z_{j,m,c}^{P_{j,m}},\qquad
 \phi(\tau_{j,m,c})=z_{j,m,c};
\]
tangents with $m=0$ are excluded. A \emph{good} monomial has label exponents $<P_{j,m}$, the prescribed degrees after $\phi$, and tangent degree at $(j,m)$ at most $d_{L,j}:=\sum_{x\in L_j}d_x$.

The chain sets $S_{j,i}$ consist of a nested subset of $L_j$ and all carries from $j-1$. The link sets are
$Y_{j,m}=L_j\cup\{u_{j-1,*}\}\cup\{\tau_{j,m'}:m'>m\}$.
The input $F$ has $b(F)=0$, prescribed multidegree, chain orders $\theta_{j,i}$ in every digit component, and link orders $\theta^L_{j,m}$ before splitting. Let $\Gamma$ multiply coefficient $\alpha$ by the prescribed unit $\gamma(\alpha)$, whose two signs admit integral binomial-polynomial approximations modulo $p^h$ of degree $\le6h$.

Set
\[
 \kappa=\sum_x r\bigl(\log_p(d_x+r)+1\bigr),\quad
 e=r^2\bigl(\log_p(2\sum_xd_x+r)+1\bigr),\quad
 D_m=e\sum_j n_j,\quad D_l=eJ(r-1).
\]
Here $\kappa$ is the source's \texttt{kp}. Assume $c\ge1$,
\[
 c>4(D_m+D_l+3\kappa),\qquad
 \frac{r d_{L,j}}{P_{j,m}}+26rc\le s_{j+1,i}.
 \tag{PG.1}
\]
Write $h=\lceil c\rceil\le2c$. Thus the last inequality holds with $13rh$ in place of $26rc$. \textsf{[PG1]}

\subsection{Constructing a detector and imposing the chains}

The weighted lattice $\{b\ge0\}$ is stable under divided powers: a root operator of order $q$ changes the weighted coefficient valuation by
\[
 \tfrac12v_p\binom{\alpha_b}{q}
 +\tfrac12v_p\binom{\alpha_a+q}{q}\ge0.
\]
The integral-projector lemma therefore gives each order-cutoff projector a loss at most $e$ in $b$. Projectors for nested or disjoint letter sets commute and are Fock self-adjoint. Choose $\alpha$ attaining $b(F)=0$, and project $x^\alpha$ onto all original link conditions. The pairing with $F$ is unchanged. Approximate $\Gamma^{-1}$ modulo $p^h$, then clear its coefficient denominators by $p^{\lceil n_0+D_l-v_p(\alpha!)/2\rceil}$. This gives an integral detector $H$ with link loss $6h$ and
\[
 v\langle F,\Gamma H\rangle\le n_0+D_l+1.
 \tag{PG.2}
\]
The signed-path basis lifts $H$ to an integral good polynomial $Z$ with $\phi Z=H$ and the same link bounds. Explicitly, in each path-column product move the frame letters $m>0$ to tangents exactly when that column contains a base letter; each tangent then occurs at most $d_{L,j}$ times. Split the remaining labels into residue and carry exponents. The column order bounds prove that this preserves the required links. \textsf{[PG3--PG5]}

Let $\Pi$ be the product of the chain projectors at levels $\theta-s$. It preserves goodness and the link conditions. On a good monomial its coefficient denominators cost at most $D_m+\kappa/2$. Factorial splitting supplies a unit function
$\gamma_*(\zeta)=\gamma(A)\gamma_{\rm dig}(\zeta)$, where $\phi(x^\zeta)=x^{A+M}$ and $M$ is the tangent exponent. Choose integral diagonal approximations $g_\pm\equiv\gamma_*^{\pm1}\pmod{p^h}$, each of order loss at most $12h$, and set
\[
 \Omega=\operatorname{diag}(g_-)\Pi\operatorname{diag}(g_+),
 \qquad G=p^{N_2}\Omega Z,\qquad N_2=\lceil D_m+\kappa/2\rceil.
 \tag{PG.3}
\]
Then $G$ is integral and good, with chain loss $12h$ and total link loss $30h$. \textsf{[PG6,PG8]}

\subsection{Why the constrained projection preserves detection}

The decisive derivative estimate is elementary. For $P=p^g$,
\[
 a_l=[z^M](1+z)^\rho\bigl((1+z)^P-1\bigr)^l,
 \quad
 \binom{\rho+P\beta}{M}=\sum_l a_l\binom\beta l,
 \quad p^{\max(0,l-\lfloor M/P\rfloor)}\mid a_l.
\]
Indeed $(1+z)^P-1=z^P+pQ(z)$: a term contributing degree $M$ can use at most $\lfloor M/P\rfloor$ copies of $z^P$. Tensoring and discarding coefficients divisible by $p^h$ replaces divided differentiation of order $M$ by carry differentiation of order at most $|M|/P+h$. Together with the two $6h$ unit corrections, (PG.1) proves that every digit component of $\Gamma^{-1}\partial^{[M]}\Gamma F$ satisfies the weakened chains, up to coefficients of valuation $\ge h-n_0$. \textsf{[PG2,PG7]}

Put $\Lambda(W)=\langle F,\Gamma\phi W\rangle$. The factorial identity on each fixed label/tangent block rewrites $\Lambda$ as a constant times a Fock pairing with that conjugated derivative. Its chain-compatible part is fixed by $\Pi$, hence cancels in $\Lambda(\Omega Z)-\Lambda(Z)$. The explicit error calculation gives
\[
 v\bigl(\Lambda(\Omega Z)-\Lambda(Z)\bigr)
 \ge h+n_0-2\kappa-D_m>n_0+D_l+1.
 \tag{PG.4}
\]
Thus (PG.2) survives. Since $\kappa\ge1$, integrality and factorial bounds give
\[
 b(\phi G)\ge n_0-\kappa,\qquad
 v\Lambda(G)\le n_0+D_m+D_l+3\kappa.
 \tag{PG.5}
\]
Finally $12h\le24c$ and $30h\le60c$. These are exactly the required orders and Gram bound, with $E=D_m+D_l+2\kappa$ and \emph{$R=0$}. The report supplies the block cancellation in (PG.4); this is a construction of the constrained test, not an assumption that a suitable nonzero pairing exists. \textsf{[PG9--PG11]}
